\subsection{Cross-references}
Everything is referenced through \verb|cleveref|, so the leading word is generated rather than typed.
Write \verb|\cref{tab:simple}| and get \cref{tab:simple}, or \verb|\Cref{fig:single}| at the start of a sentence and get \Cref{fig:single}.
Several references collapse automatically, so \verb|\cref{tab:simple,tab:grouped,tab:wrapping}| gives \cref{tab:simple,tab:grouped,tab:wrapping}.
Never hard-code the word ``Table'' or ``Figure'' in front of a reference, because it defeats the whole mechanism.

Label prefixes keep the reference space legible.
Use \verb|fig:|, \verb|tab:|, \verb|eq:|, \verb|sec:|, and \verb|alg:|.

\subsection{Citations}
Citations are numeric and sorted, through \verb|natbib|.
Use \verb|\citep| for a parenthetical citation, as in \citep{example2026}, and \verb|\citet| for a textual one, as in \citet{example2026}.
Several keys in one command compress into a range, so \verb|\citep{example2026,example2025}| gives \citep{example2026,example2025}.
Page numbers go in the optional argument, as in \citep[p.~14]{example2026}.

\subsection{Equations}
A numbered display uses \verb|equation|, and stays inside the draft-mode line numbering.

\begin{equation}\label{eq:mutual}
  I(A;B) = H(A) + H(B) - H(A,B)
\end{equation}

Use \verb|align| for several related lines, aligned on the relation symbol.
Mark any line that does not need a number with \verb|\nonumber|.

\begin{align}
  R &\geq q\left[1 - H_2(e_\mathrm{b}) - f H_2(e_\mathrm{p})\right] \nonumber \\
    &\geq q\left[1 - 2H_2(e)\right]
  \label{eq:keyrate}
\end{align}

An unnumbered display uses \verb|\[ ... \]|, and a case distinction uses \verb|cases|.

\[
  \rho = \begin{cases}
    \ket{0}\bra{0} & \text{with probability } p, \\
    \ket{1}\bra{1} & \text{otherwise.}
  \end{cases}
\]

\Cref{eq:mutual,eq:keyrate} are referenced the same way as everything else.

\subsection{Lists}
\verb|enumitem| controls the spacing.
The \verb|noitemsep| key removes the gaps between items, which matters in a two-column paper where vertical space is scarce.

\begin{itemize}[noitemsep]
  \item A tight bulleted item.
  \item A second one.
\end{itemize}

\begin{enumerate}[label=(\roman*),noitemsep]
  \item A numbered item with a custom label format.
  \item A second one.
\end{enumerate}

\begin{description}[noitemsep]
  \item[Concept] The phase in which the problem space is characterised.
  \item[Definition] The phase in which requirements and architecture are settled.
\end{description}

\subsection{Other fragments}
A footnote looks like this.\footnote{Footnotes sit at the foot of the column in the two-column model.}
Inline verbatim uses \verb|\verb|, and a short quotation uses the \verb|quote| environment.

\begin{quote}
  A displayed quotation, indented on both sides and set without quotation marks.
\end{quote}
